% !TeX program = xelatex
% !TeX encoding = UTF-8
\documentclass{JXUSTmodeling}
\geometry{reset,a4paper,left=2.5cm,right=2.5cm,top=2.5cm,bottom=2.5cm}
\usepackage{float}
\graphicspath{{figures/}{./}}
\hypersetup{pdftitle={基于相容区域与覆盖认证的无线电干扰源定位清除},pdfauthor={},pdfsubject={},pdfkeywords={有界误差,半平面交,最坏直径,方向覆盖}}
\setlength{\emergencystretch}{2em}
\setlength{\tabcolsep}{5pt}
\labelformat{figure}{#1}
\labelformat{table}{#1}
\labelformat{equation}{#1}
\biaoti{基于相容区域与覆盖认证的\\无线电干扰源快速定位与清除}
\keyword{有界测向误差；半平面交；最坏直径；定向覆盖；滚动路径规划}
\begin{document}
\begin{abstract}
针对机器狗在有限接收距离及有界测向误差下的无线电干扰源定位清除问题，本文以相容位置集合为基础，将交会定位、主动选点与覆盖搜索统一为逐步缩小不确定区域的决策过程，在保证清除完整性的前提下减少移动及检测开销。

针对问题一，将测向扇区转化为半平面交，分别识别空集、无界和退化情形，对有界凸多边形采用旋转卡壳计算直径。通过可实现的等边三角形观测反例说明：以区域直径为直径的圆未必覆盖全域。进一步以最小覆盖圆半径建立20米光学定位的充分必要判据。

针对问题二，利用首次收信隐含的接收半径下界，推导四圆盘交集安全域；以最坏相容读数下的区域直径为目标，结合解析条带上界、目标圆条件化及读数网格半步增宽进行分层选点。在两个自建场景中，同口径直径上界分别由115.693米和54.024米降至112.329米和15.769米，结果为有限候选中的优选。

针对问题三，根据原点收信情况选择六点内环或九点外环覆盖，将全向阴性约束、有限前瞻选点与搜索路径联合调度，并在细长区域比较窄带光学清除与射频定位成本。完整光学网格仍作为有限兜底。在512例新留出案例中全部清除并完成确认，平均226.8秒/源，较同例原默认减少16.8\%。

针对问题四，将任意定向朝向可见性转化为邻近检测点的凸包条件，利用有限Voronoi圆缺极值检验认证22点布局，结合双阴性截断、共享基线与滚动路径优化。以假想观测收缩量筛选顺路补测，并在清除失败后保留剩余碎片并集，按覆盖面积与动作成本选择有界光学修复。在512例新留出案例中全部清除并完成确认，平均441.7秒/源，较同例原默认减少9.0\%。

本文区分数学保证、自建实验和可核对的接口记录；新版结果以冻结后的独立样本检验，旧版平台演练不作为新版正式成绩。模型在题设条件与足够执行预算下具有完整性依据，但有限搜索和路径启发式不保证全局最优。
\end{abstract}

\section{问题重述}\label{sec:1}
\subsection{问题背景}\label{sub:1.1}
无线电干扰源的精准定位与清除是无线电频谱管理的重要任务。传统作业依赖技术人员携带测向设备徒步排查，在危险环境中存在安全风险，定位效率也受到限制。为实现自动化作业，现拟利用搭载测向机及光学定位、清除装置的机器狗，通过多点方位观测确定干扰源位置。

在完成初步监测和区域性排查后，精准定位仍受到测向误差与有限接收距离的制约，单点观测难以唯一确定干扰源位置。对于数量未知的多目标场景，频繁切换频道、重复检测及不合理移动还会增加作业时间。因此，需要统筹区域搜索、检测点选择与定位清除过程，在保证清除完整性的同时提高整体作业效率。

\subsection{问题提出}\label{sub:1.2}
已知干扰源位于半径为1800米的圆形目标区域内，以圆心为原点，正东、正北分别为横、纵轴正方向，示向度由横轴正向逆时针计量。误差不超过$1^\circ$且同一地点误差固定，有效接收半径在1000至1500米之间、具体值未知。各源频道互异，取自20个频道。围绕自动定位与清除，需解决以下四个问题。

\textbf{问题一：}根据若干检测点坐标及其对同一干扰源的示向度，设计交会定位区域直径的计算算法，并判断以该直径为直径的圆能否覆盖整个区域。

\textbf{问题二：}对于全向干扰源，在已知一个检测点及其示向度的条件下，制定第二检测点的选择策略，并给出能够获得较好定位效果的候选区域。

\textbf{问题三：}区域内有10至16个全向源，具体数量未知。机器狗由原点出发、初始频道为1，制定自动搜索、定位与清除策略，在确保全部清除的前提下尽量缩短总任务时间，并通过模拟器演练评价清除比例和平均定位清除时间。

\textbf{问题四：}区域内同时存在全向与定向源，总数仍为10至16个，定向源数量与方向均未知。定向源只在其发射方向两侧各$90^\circ$内辐射，其余条件与问题三相同。制定适应混合场景的检测清除策略，并通过模拟器检验。

问题三、四均要求在充分演练后分别完成三次正式测试，按题设记录案例编码、清除数量、平均定位清除时间及程序运行时间，并将原名加密日志纳入支撑材料。

\section{问题分析}\label{sec:2}
\subsection{问题一}\label{sub:2.1}
考虑测向误差后，单次观测将源位置约束在示向线两侧各$1^\circ$的角扇区内，多次观测的交集即为交会定位区域。可将各扇区转化为两个半平面约束，先判定交集状态，再利用凸性把直径转化为最远顶点对距离。覆盖性须独立判断：最大跨度不能直接确定覆盖全域所需的圆半径，应结合反例与最小覆盖圆作答。

\subsection{问题二}\label{sub:2.2}
第二点选择需兼顾再次收信与交会几何。两次测向方向接近共线会放大位置歧义，而任意横向移动又可能超出未知接收范围。因此，先利用首次收信信息构造保守安全域，再结合目标圆约束，以第二次观测后区域直径的最坏上界评价选点。用分层搜索给出优选点及近优候选区域，并统一评分精度，区分上界变紧与真实误差改善。

\subsection{问题三}\label{sub:2.3}
源数量未知，单纯追逐已发现目标无法证明任务完成，必须保留覆盖全域的搜索依据。全向源的接收半径下界可将发现过程转化为有限圆盘覆盖。由于移动与逐频道检测均有成本，应将发现新源和清除已知源联合调度，并利用途中停点补充观测。对于几何交会效果较差的区域，采用有限光学覆盖保证已知源最终能够被清除。

\subsection{问题四}\label{sub:2.4}
定向源失收可能源于距离过远，也可能源于处在发射背面，故不能直接照搬全向覆盖或将阴性点周围圆盘删除。应构造对所有朝向均有效的检测点集，并为连续目标域提供覆盖证明。在定位阶段，仅在接收距离已获保证时，结合阳性点与成对阴性点提取几何约束；随后通过历史观测复用与滚动路径改进降低额外成本。

\section{模型的假设}\label{sec:3}
在题设基础上，采用以下假设与计算约定。

\textbf{（1）平面静态场景。}忽略高度差，在二维平面内建模；一次任务期间源位置、频道、发射方向及接收半径保持不变。

\textbf{（2）点位与运动准确。}机器狗能够准确获知并到达指定检测点，路径可无障碍直达；忽略自身定位误差、转向和加减速额外耗时，直线速度为5米/秒。

\textbf{（3）固定有界误差。}不设测向误差概率分布或独立性，不利用同点重复观测平均缩小误差界。前两问按$1^\circ$建模；第三、四问将两位小数读数的舍入误差并入，取$1.005^\circ$。

\textbf{（4）未知参数保守处理。}位置、半径和方向仅受题设与已有反馈约束，策略不访问源真值。第二问候选点对全部相容状态保证收信；1800米圆仅约束源的位置。第三、四问假想读数、半径先验与面积权重只用于动作排序，不作为删去候选源位置或认定完成的依据。

\textbf{（5）近源与清除。}信号覆盖内距源不超过5米时按近源反馈处理；光学定位及清除只取决于20米距离条件，不受发射方向限制。检测耗时5秒、换频1秒；清除成功计光学3秒与激光2秒，失败仅计光学3秒。

\textbf{（6）反馈与预算。}状态仅在接口确认动作被接受后更新。数学完整性结论以动作正常执行、误差界成立及预算足够为前提；拒绝、连接异常或预算耗尽不视为成功完成。本地仿真的误差模式和源分布只用于测试，不当作官方数据生成规律。

\section{符号说明}\label{sec:4}
\begin{table}[htbp]
\centering
\begin{tabularx}{\textwidth}{cXc}
\toprule
符号 & 含义 & 单位\\
\midrule
$\Omega$ & 半径为1800米的目标圆域 & --\\
$G,g$ & 干扰源的真实位置（仅用于数学分析与事后评价） & 米\\
$S_i,p_i$ & 第$i$个检测点位置 & 米\\
$\theta_i,\varepsilon$ & 示向度及采用的测向误差界 & 度\\
$R$ & 未知有效接收半径，$1000\le R\le1500$ & 米\\
$P,U$ & 与已有观测相容的位置集合或其保守外包 & --\\
$D(P),R_*(P)$ & 区域直径、最小覆盖圆半径 & 米\\
$q=(a,b)$ & 第二检测点在首示向局部坐标中的位置 & 米\\
$J(q),A(q),B(q),C(q)$ & 最坏直径目标及其解析、几何、组合上界 & 米\\
$V_3,V_4$ & 全向与混合场景的有限覆盖检测点集 & --\\
$N_c,N$ & 已清除数、真实总源数（$N$不输入策略） & 个\\
$T,\overline T$ & 虚拟总任务时间、每源平均定位清除时间 & 秒\\
\bottomrule
\end{tabularx}
\end{table}
其他局部符号在首次使用时说明。$B(p,r)$表示以$p$为圆心、$r$为半径的圆盘，与第二问上界函数$B(q)$按参数形式区分。

\section{模型的建立与求解}\label{sec:5}
本文以所有与题设及反馈相容的位置构成定位结果，将新增观测转化为集合约束。误差扇区的半平面表达与有界方位定位研究相一致\cite{gholami2015}。第一问处理几何评价，第二问处理单次主动选点，第三、四问进一步将集合更新与有限覆盖和动作调度结合。
